%==============================================================================
\section{Signal Processing Pipeline}
%==============================================================================

Raw ECG captured from a wearable, battery-powered node is noisy: it suffers from
baseline wander (slow drift as the patient breathes and moves), 50/60\,Hz mains
interference, muscle (EMG) noise, and occasional motion artefacts. The function
\texttt{filter\_ecg\_signal()} in \texttt{ecg/filters.py} cleans each incoming
batch before it is broadcast. It is built on NumPy, SciPy's signal toolbox
(\texttt{butter}, \texttt{sosfiltfilt}, \texttt{iirnotch}, \texttt{savgol\_filter},
\texttt{find\_peaks}, \texttt{welch}) and PyWavelets.

The default sampling rate is \texttt{fs = 250}\,Hz, matching the firmware. The
pipeline runs the eight stages shown in Figure~\ref{fig:dsp} and described below.

\begin{figure}[h]
  \centering
  \begin{tikzpicture}[node distance=5.5mm]
    \node[stage] (s0) {\textbf{0. Outlier clipping} -- robust clip to
        $\text{median} \pm 5\,\text{MAD}$};
    \node[stage, below=of s0] (s1) {\textbf{1. Baseline-wander removal} --
        segmented polynomial fit + 0.5\,Hz high-pass (Butterworth, SOS)};
    \node[stage, below=of s1] (s2) {\textbf{2. Three-band filtering} --
        wide 0.5--45\,Hz $\cdot$ QRS 8--20\,Hz $\cdot$ P/T 0.5--10\,Hz};
    \node[stage, below=of s2] (s3) {\textbf{3. Adaptive QRS weighting} --
        combine bands $0.3 / 0.5 / 0.2$ using a smoothed QRS mask};
    \node[stage, below=of s3] (s4) {\textbf{4. Powerline removal} --
        auto-detect 50/60\,Hz, cascaded notch on harmonics $1\times,2\times,3\times$};
    \node[stage, below=of s4] (s5) {\textbf{5. Wavelet denoising} --
        \texttt{sym6}, level-dependent soft thresholding};
    \node[stage, below=of s5] (s6) {\textbf{6. Adaptive smoothing} --
        SNR-based Savitzky--Golay or 35\,Hz low-pass};
    \node[stage, below=of s6] (s7) {\textbf{7. R-peak-preserving outlier replace}
        -- rolling median/MAD, peaks protected};
    \node[stage, below=of s7] (s8) {\textbf{8. Re-centre \& clip} --
        centre on 2048, clip to ADC range 0--4095};

    \foreach \a/\b in {s0/s1,s1/s2,s2/s3,s3/s4,s4/s5,s5/s6,s6/s7,s7/s8}
        \draw[flow] (\a) -- (\b);

    \node[hw, left=8mm of s0, minimum width=18mm] (in) {raw\\\scriptsize ecg[]};
    \node[fe, left=8mm of s8, minimum width=18mm] (out) {clean\\\scriptsize filtered[]};
    \draw[flow] (in) -- (s0);
    \draw[flow] (s8) -- (out);
  \end{tikzpicture}
  \caption{The eight-stage \texttt{filter\_ecg\_signal()} pipeline. Each batch of
  raw samples flows top-to-bottom and emerges as a cleaned, ADC-scaled array.}
  \label{fig:dsp}
\end{figure}

\subsection*{Stage details}
\begin{description}[leftmargin=0pt,style=nextline]
  \item[0. Robust outlier clipping] Values are clipped to
        $\text{median} \pm 5\sigma$, where $\sigma$ is estimated from the median
        absolute deviation (MAD, scaled by $1.4826$). This removes gross spikes
        without being skewed by them, unlike a mean/standard-deviation clip.
  \item[1. Baseline-wander removal] For batches longer than five seconds the
        signal is detrended with a piecewise (segmented) polynomial fit to track
        non-stationary drift; then a 0.5\,Hz Butterworth high-pass
        (second-order-sections, zero-phase \texttt{sosfiltfilt}) removes the
        remainder.
  \item[2. Three-band filtering] The signal is split into three Butterworth
        band-passes tuned to different ECG features: a \emph{wide} band
        (0.5--45\,Hz) preserving overall morphology, a \emph{QRS} band (8--20\,Hz)
        isolating the sharp ventricular spike, and a \emph{P/T} band (0.5--10\,Hz)
        capturing the slower atrial and repolarisation waves.
  \item[3. Adaptive QRS weighting] A QRS energy mask is built with a sliding
        adaptive threshold ($\text{mean} + 1.8\,\text{std}$ over 1.5\,s windows)
        and smoothed. The three bands are recombined as
        $0.3\cdot\text{wide} + 0.5\cdot\text{QRS}\!\cdot\!\text{mask}
        + 0.2\cdot\text{P/T}\!\cdot\!(1-\text{mask})$, so the QRS band dominates
        during beats and the gentler bands dominate between them.
  \item[4. Powerline removal] If the mains frequency is not given, it is
        auto-detected via Welch power-spectral-density peaks near 50 or 60\,Hz.
        Cascaded \texttt{iirnotch} filters (two Q-factors, 35 and 50) then notch
        out the fundamental and its 2nd and 3rd harmonics.
  \item[5. Wavelet denoising] A \texttt{sym6} (Symlet-6) wavelet decomposition with
        level-dependent soft thresholding suppresses broadband noise while
        preserving ECG morphology; low-frequency detail is thresholded gently and
        high-frequency detail more aggressively.
  \item[6. Adaptive smoothing] The signal-to-noise ratio is estimated; very noisy
        batches get a short ($\approx$40\,ms) Savitzky--Golay smooth, moderately
        noisy batches a 35\,Hz low-pass, and clean batches are left untouched.
  \item[7. R-peak-preserving outlier replacement] R-peaks are located with the QRS
        band and an adaptive height threshold. A rolling median/MAD test then
        replaces residual outliers, but explicitly \emph{protects} the detected
        R-peak regions so the diagnostically important spikes are not flattened.
  \item[8. Re-centre and clip] Finally the signal is centred on 2048 (the middle of
        the 12-bit ADC range) and clipped to integer codes $0$--$4095$, matching
        the dashboard's fixed $y$-axis.
\end{description}

\paragraph{Output contract.} The function always returns a list of integers in
$[0,4095]$. Short batches (fewer than 20 samples) are returned unchanged, and any
stage that fails (for example wavelet reconstruction on an awkward length) falls
back gracefully so a single bad batch never breaks the stream.
